% Contribution for the continuing Polylogarithms report.
% Requires amsmath, amssymb, amsthm, booktabs, hyperref.
% Citation keys: Ljunggren1960, Tverberg1960, BoydMartinThom2015,
% AartsFokkinkKruijtzer2001.

\section{Rigidity of complementary powers}
\label{sec:complementary-rigidity}

The earlier complementary-power ladder starts from a relation
$r^a+r^b=1$, where $0<r<1$ and $a,b$ are positive integers.
That construction raises a concrete completeness question: can one base
have many different pairs of complementary powers?  We settle this
question for pure powers, including arbitrary common divisors in the
exponents.  The decisive arithmetic input is a classical irreducibility
theorem for trinomials.  Our conclusions concern this precisely specified
class of seeds; they do not classify general multiplicative relations
among $1-r^k$, or bound the maximum weight of polylogarithm ladders.

For $0<r<1$, define
\[
\mathcal C(r)=\{(a,b)\in\mathbb Z_{>0}^2:
                 a\le b,\ r^a+r^b=1\}.
\]
Let $\omega$ be the unique root in $(0,1)$ of
$X^3+X^2-1$.  Thus $\omega$ is the reciprocal of the plastic number.

\begin{theorem}[Complete classification of collisions]
\label{thm:complement-collisions}
For every real $0<r<1$, the set $\mathcal C(r)$ has at most two elements.
It has two elements if and only if $r^d=\omega$ for some positive
integer $d$, in which case
\[
\mathcal C(r)=\{(d,5d),(2d,3d)\}.
\]
If $\mathcal C(r)$ contains $(a,a)$, then
$\mathcal C(r)=\{(a,a)\}$ and $r=2^{-1/a}$.
\end{theorem}

The theorem supplies an exact answer to the seed-classification question
posed after the earlier cubic and quartic ladder certificates.  It is
proved below as a consequence of the Ljunggren--Tverberg theorem, with
all additional arguments included.  We make no claim that this
consequence is globally new in the literature.  The related theorem of
Aarts, Fokkink, and Kruijtzer classifies the two morphic numbers, by
requiring both $p+1$ and $p-1$ to be powers of $p$; their proof also uses
trinomial irreducibility \cite{AartsFokkinkKruijtzer2001}.  Our formulation
permits arbitrary positive exponent pairs and explicitly handles power
substitutions.

\subsection{The irreducibility input and the entire cyclotomic factor}

For $1\le a<b$, put $d=\gcd(a,b)$, $A=a/d$, $B=b/d$, and
\[
F_{a,b}(X)=X^b+X^a-1.
\]
The following precise specialization is the arithmetic input.

\begin{theorem}[Ljunggren--Tverberg, specialized]
\label{thm:trinomial-input}
The polynomial $F_{a,b}$ has exactly one irreducible factor over
$\mathbb Q$ whose roots are not roots of unity.  Its other factor is
\[
C_{a,b}(X)=
\begin{cases}
 X^{2d}-X^d+1,&(A,B)\equiv(1,5)\text{ or }(5,1)\pmod6,\\
 1,&\text{otherwise}.
\end{cases}
\]
In particular $H_{a,b}=F_{a,b}/C_{a,b}$ is irreducible, monic,
and has constant term $-1$.  Moreover
$H_{a,b}(X)=H_{A,B}(X^d)$.
\end{theorem}

This is the trinomial result of Ljunggren and Tverberg
\cite{Ljunggren1960,Tverberg1960}; a modern explicit formulation,
including the assertion that $H_{A,B}(X^d)$ remains irreducible, is
Boyd--Martin--Thom, Lemma~4.1(a),(b) \cite{BoydMartinThom2015}.
Only the trinomial theorem is used.  Its separate quadrinomial analogue
has a more delicate history and plays no role here.

For clarity, the displayed cyclotomic factor can be obtained directly,
without relying on a table of sign cases.  Suppose $|z|=1$ and
$z^a+z^b=1$.  The two unit vectors must be the two numbers
$e^{i\pi/3}$ and $e^{-i\pi/3}$, in either order.  If $w=z^d$, then
$w^A$ and $w^B$ are these primitive sixth roots.  A Bezout identity for
$A,B$ implies $w^6=1$, and $w^A$ of order six then implies that $w$
itself has order six.  Consequently $A,B$ have opposite residues
$1,5$ modulo six.  Conversely those residues make every root of
$X^{2d}-X^d+1$ a root of $F_{a,b}$.  These roots are simple: a common
zero of $F_{a,b}$ and its derivative on the unit circle would give
$a|z^a|=b|z^b|$, hence $a=b$, a contradiction.  Thus the formula for
$C_{a,b}$ accounts for the \emph{entire} cyclotomic factor.  The
remaining assertion, irreducibility of $H_{a,b}$, is the cited theorem.

\subsection{A conjugate-rotation invariant}

\begin{lemma}[The common divisor is intrinsic]
\label{lem:intrinsic-gcd}
Suppose $r^a+r^b=r^c+r^e=1$, where $0<r<1$,
$1\le a<b$, and $1\le c<e$.  Then
\[
\gcd(a,b)=\gcd(c,e).
\]
\end{lemma}

\begin{proof}
Set $d=\gcd(a,b)$.  Since $0<r<1$, it is not a root of unity, so
its minimal polynomial is $H_{a,b}$.  Theorem~\ref{thm:trinomial-input}
shows that this polynomial is a polynomial in $X^d$.  If $\zeta$ is a
primitive $d$th root of unity, $r\zeta$ is therefore a conjugate of $r$.
The second equation also holds at this conjugate:
\[
r^c\zeta^c+r^e\zeta^e=1.
\]
Taking absolute values gives equality in the triangle inequality,
because $r^c+r^e=1$.  Both summands must have the argument of their
positive real sum.  Thus $\zeta^c=\zeta^e=1$, proving
$d\mid\gcd(c,e)$.  Interchanging the pairs proves the reverse
divisibility.
\end{proof}

This step prevents a gap that would arise from normalizing the two
pairs independently.  The cyclotomic factor can have degree $2d$,
so a degree comparison before proving that the common divisors agree
would not suffice.

\subsection{The primitive coefficient comparison}

\begin{proof}[Proof of Theorem~\ref{thm:complement-collisions}]
First suppose an equal pair occurs.  Then $r^a=1/2$, so $r$ cannot
be an algebraic integer: otherwise the rational number $r^a=1/2$
would be an algebraic integer.  Any unequal pair would make $r$ a
root of the monic integer polynomial $F_{c,e}$, a contradiction.
There is only one equal pair because the powers of $r$ strictly
decrease.

Now suppose two distinct unequal pairs occur.  By
Lemma~\ref{lem:intrinsic-gcd} they have a common gcd $d$.
Replace $r$ by $r^d$ and divide all four exponents by $d$.
Both pairs are now primitive.  Their trinomials have the same
irreducible noncyclotomic factor, and their cyclotomic factors belong
to $\{1,\Phi_6\}$, where $\Phi_6=X^2-X+1$.
If the two cyclotomic factors agree, the trinomials agree, contrary
to distinctness.  After interchanging them if necessary, there must
therefore be a polynomial identity
\begin{equation}
(X^2-X+1)(X^n+X^m-1)=X^{n+2}+X^k-1,
\quad 1\le m<n,\quad 1\le k<n+2.
\label{eq:primitive-collision-factor}
\end{equation}
The coefficient of $X^{n+1}$ on the left is $-1$ unless
$m=n-1$, in which case it is zero.  The right-hand coefficient is
zero or one.  Hence $m=n-1$, and the identity reduces to
\[
X^{n+2}+X^{n-1}-X^2+X-1=X^{n+2}+X^k-1.
\]
For $n=2$ the remaining polynomial is $2X-X^2$, which is not a
monomial.  For $n>3$ its coefficient at $X^2$ is $-1$.
Thus $n=3$, which gives $m=2$ and $k=1$.
The common factor is $X^3+X^2-1$, and the two normalized pairs are
$(1,5)$ and $(2,3)$.  Conversely,
\[
X^5+X-1=(X^2-X+1)(X^3+X^2-1)
\]
shows that these two pairs do occur at $\omega$.
Finally, a third pair would have to be the same exceptional companion
of either existing pair, so no third pair occurs.
\end{proof}

\subsection{All signed pure-power unit equations}

The functional equations for polylogarithms use signs and inverses as
well as positive powers.  The preceding classification automatically
controls all such complementary arguments.

\begin{corollary}[Sharp signed-unit count]
\label{cor:signed-unit-count}
Let $\Gamma_r=\{\pm r^k:k\in\mathbb Z\}$ and
\[
\mathcal U(r)=\{x\in\Gamma_r:1-x\in\Gamma_r\}.
\]
Then $|\mathcal U(r)|$ is one of $0,3,6,12$.
For every unequal pair $(a,b)\in\mathcal C(r)$, its contribution is
the six-element set
\begin{equation}
\left\{r^a,r^b,r^{-a},r^{-b},-r^{b-a},-r^{a-b}\right\}.
\label{eq:six-unit-orbit}
\end{equation}
These sets are disjoint for different unordered pairs.
An equal pair contributes precisely $\{1/2,2,-1\}$.
The case $|\mathcal U(r)|=12$ occurs exactly when $r^d=\omega$
for some positive integer $d$.
\end{corollary}

\begin{proof}
The transformations $x\mapsto1-x$ and $x\mapsto1/x$ preserve
$\mathcal U(r)$.  For example,
$1-1/x=-(1-x)/x\in\Gamma_r$ whenever $x,1-x\in\Gamma_r$.
Their orbit consists of
\[
x,\quad1-x,\quad x^{-1},\quad(1-x)^{-1},
\quad x/(x-1),\quad(x-1)/x.
\]
Every real orbit has a representative in $(0,1)$: use $1/x$ if
$x>1$ and $1/(1-x)$ if $x<0$.  Such a representative and its
complement are $r^a,r^b$ with $a,b>0$.  The two members of a
nondegenerate real orbit that lie in $(0,1)$ are exactly $x$ and
$1-x$, so its unordered positive pair is unique.  Substituting
$x=r^a$, $1-x=r^b$ gives \eqref{eq:six-unit-orbit}.
All six entries are distinct unless $a=b$, when they reduce to
the stated three values.  The cardinalities now follow from
Theorem~\ref{thm:complement-collisions}.
\end{proof}

All four possibilities occur: $r=1/3$, $1/2$, the reciprocal golden
ratio, and $\omega$ give respectively $0,3,6,12$.
The conclusion classifies \emph{pure signed-power complements}.
It does not say that there are at most twelve useful polylogarithm
arguments: products of units $1-r^k$ can introduce many more arguments.

\begin{corollary}[Rational exponents]
\label{cor:rational-complements}
The same collision classification holds for positive rational
exponents, with the exceptional common scale $d$ allowed to be a
positive rational number.  In particular there are at most two
unordered positive rational pairs, and
\[
\#\{x\in\{\pm r^q:q\in\mathbb Q\}:1-x\in
                          \{\pm r^q:q\in\mathbb Q\}\}\le12.
\]
\end{corollary}

\begin{proof}
Given two pairs, choose a common denominator $N$ and apply
Theorem~\ref{thm:complement-collisions} to $r^{1/N}$.
A third rational pair could be included in the same denominator
and would contradict the integer theorem.  The orbit proof is
unchanged.
\end{proof}

\subsection{A provably complete bounded-degree search}

\begin{corollary}[Degree and an exhaustive exponent bound]
\label{cor:degree-search}
If $r^a+r^b=1$ with $1\le a<b$, then
\[
[\mathbb Q(r):\mathbb Q]=
\begin{cases}
b-2d,&(a/d,b/d)\equiv(1,5)\text{ or }(5,1)\pmod6,\\
b,&\text{otherwise},
\end{cases}
\qquad d=\gcd(a,b).
\]
Consequently all such bases of degree at most $D$ are found by
enumerating
\[
1\le a<b\le\left\lfloor\frac{5D}{3}\right\rfloor
\]
and retaining those whose displayed degree is at most $D$.
For primitive pairs it suffices to enumerate $b\le D+2$.
The only repeated bases in either enumeration are the classified
plastic pairs.  Equal pairs add the separate bases $2^{-1/a}$,
with $a\le D$.
\end{corollary}

\begin{proof}
The degree formula is Theorem~\ref{thm:trinomial-input}.  In the
exceptional case $b/d\ge5$, so $b-2d\ge3d$.  If the degree is
at most $D$, then $d\le D/3$ and
$b=(b-2d)+2d\le5D/3$.  The nonexceptional case is immediate.
For primitive pairs the possible degree loss is only two.
The equal-pair polynomial $2X^a-1$ has degree $a$ and is
irreducible: its reciprocal is, up to a nonzero scalar,
$X^a-2$, which is Eisenstein at $2$.
\end{proof}

This bound is sharp for infinitely many $D$: the pair $(d,5d)$
has degree $3d$ and largest exponent $5d$.
Unlike a search over bounded polynomial coefficients, the algorithm
has an explicit completeness theorem for its stated class.  No
numerical equality test is required.  A base is stored by the exact
irreducible polynomial $H_{a,b}$ together with its unique root in
$(0,1)$, and duplicate detection is exact polynomial comparison.

For illustration, the complete unequal-pair list through degree four is
\[
\begin{array}{ccl}
\text{degree}&\text{complement pair(s)}&\text{minimal polynomial}\\\hline
2&(1,2)&X^2+X-1\\
3&(1,3)&X^3+X-1\\
3&(1,5),(2,3)&X^3+X^2-1\\
4&(1,4)&X^4+X-1\\
4&(2,4)&X^4+X^2-1\\
4&(3,4)&X^4+X^3-1
\end{array}
\]
The degree-four entry with pair $(2,4)$ is exactly a square root of
the reciprocal golden ratio.  Its presence illustrates why normalization
of exponents is necessary when counting bases rather than primitive
seed types.

The exact enumeration in the accompanying certificate gives:
\[
\begin{array}{c|rrrrrrrrrrrr}
D&1&2&3&4&5&6&7&8&9&10&11&12\\\hline
\text{unequal-pair bases}&0&1&3&6&10&15&20&27&37&46&56&67\\
\text{primitive bases}&0&1&3&5&9&11&16&20&28&32&42&46\\
\text{all bases}&1&3&6&10&15&21&27&35&46&56&67&79
\end{array}
\]
Each row counts bases of degree at most $D$.  A primitive base has
an unequal complementary pair of gcd one; the intrinsic-gcd lemma
makes this definition independent of the choice of pair.
The last row also includes the $D$ equal-pair bases $2^{-1/a}$,
$1\le a\le D$.

\subsection{An exact counting formula}

The classification also counts bases without constructing or factoring
their minimal polynomials.  For $n\ge1$ define
\[
e(n)=\#\{a:1\le a<n,\ \gcd(a,n)=1,
            \ (a,n)\equiv(1,5)\text{ or }(5,1)\pmod6\}.
\]
In particular $e(n)=0$ unless $n\equiv1$ or $5\pmod6$,
and $e(1)=0$.

\begin{corollary}[Counting all bases by algebraic degree]
\label{cor:base-count}
Let $N(D)$ be the number of real $0<r<1$ with
$[\mathbb Q(r):\mathbb Q]\le D$ and $\mathcal C(r)\ne\varnothing$.
For every positive integer $D$,
\begin{align}
N(D)&=\frac{D(D+1)}2-\left\lfloor\frac D3\right\rfloor\notag\\
&\quad+\sum_{d=1}^{\lfloor D/3\rfloor}
\left\{e\!\left(\left\lfloor\frac Dd\right\rfloor+1\right)
       +e\!\left(\left\lfloor\frac Dd\right\rfloor+2\right)\right\}.
\label{eq:exact-base-count}
\end{align}
In particular,
\[
N(D)=\frac12D^2+O(D\log D),\qquad D\longrightarrow\infty.
\]
\end{corollary}

\begin{proof}
There are $D(D-1)/2$ unequal pairs with largest exponent at most
$D$, and their associated bases all have degree at most $D$.
A pair with largest exponent greater than $D$ can still have degree
at most $D$ only in the exceptional cyclotomic case.  Write that
pair $(dA,dB)$, where $\gcd(A,B)=1$ and the residues are opposite
$1,5$ modulo six.  The two inequalities are
\[
dB>D,\qquad d(B-2)\le D.
\]
Equivalently,
$B\in\{\lfloor D/d\rfloor+1,\lfloor D/d\rfloor+2\}$.
Since the exceptional normalized degree $B-2$ is at least three,
only $d\le D/3$ can occur.  This gives the sum in
\eqref{eq:exact-base-count}.
Every base counted twice is $\omega^{1/d}$, of degree $3d$;
there are exactly $\lfloor D/3\rfloor$ such duplicates.
Finally the $D$ equal-pair bases are distinct from all unequal-pair
bases and from each other.  This proves the exact formula.

For the asymptotic statement, use $0\le e(n)\le n$ to bound the sum
by
\[
\sum_{d\le D/3}\left(2\left\lfloor\frac Dd\right\rfloor+3\right)
\le2D\sum_{d\le D/3}\frac1d+D.
\]
The harmonic sum is $O(\log D)$, proving the claim.
\end{proof}

For an alternative exact implementation, when $n>1$ is coprime to
six, Mobius inversion gives
\begin{equation}
e(n)=\sum_{d\mid n}\mu(d)
       \left\lfloor\frac{n/d+1}{6}\right\rfloor.
\label{eq:exceptional-count-mobius}
\end{equation}
Indeed, after writing $a=dj$, the condition
$a\equiv-n\pmod6$ becomes $j\equiv-(n/d)\pmod6$ because every
unit modulo six is its own inverse.  For $t$ coprime to six,
the number of integers $1\le j<t$ in the residue class $-t$
is $\lfloor(t+1)/6\rfloor$.  Summing with $\mu(d)$ imposes
$\gcd(a,n)=1$.

% Citation keys: ZagierPolylogs1991, NakamuraShiraishi2025.

\section{Exact proofs and log-free ladders at the plastic base}
\label{sec:plastic-proofs}

The exceptional base in Theorem~\ref{thm:complement-collisions}
has two complementary pairs.  This additional structure produces
short proofs of two previously numerically supported formulas in
the manuscript, as well as log-free identities for $\zeta(2)$ and
$\zeta(3)$.

Throughout this section, $0<q<1$, $q^2+q^3=1$, and
$\lambda=\log q$.  The factorization in the preceding proof gives
\begin{equation}
1-q=q^5,\qquad1-q^2=q^3,\qquad
1-q^3=q^2,\qquad1-q^5=q.
\label{eq:plastic-complements}
\end{equation}
All logarithms used in the final identities are real.

\subsection{The dilogarithm formula: one Rogers pentagon}

Write
\[
R(x)=\operatorname{Li}_2(x)+\frac12\log x\log(1-x),
\qquad 0<x<1.
\]
The Rogers reflection and pentagon equations are
\begin{align}
R(x)+R(1-x)&=\zeta(2),\label{eq:plastic-R-reflect}\\
R(x)+R(y)&=R(xy)
 +R\!\left(\frac{x(1-y)}{1-xy}\right)
 +R\!\left(\frac{y(1-x)}{1-xy}\right).
\label{eq:plastic-R-pentagon}
\end{align}
They hold for $0<x,y<1$ and follow from the classical dilogarithm
functional equations; alternatively one can differentiate in $x$
and evaluate at $x=0$.  This real-domain formulation avoids branch
choices \cite{ZagierPolylogs1991}.

\begin{theorem}[The manuscript's plastic dilogarithm formula]
\label{thm:plastic-dilog}
\begin{equation}
\frac12\operatorname{Li}_2(q^2)-\operatorname{Li}_2(q^5)
=\lambda^2.
\label{eq:plastic-dilog-proved}
\end{equation}
\end{theorem}

\begin{proof}
Set $R_k=R(q^k)$.  At $(x,y)=(q,q^2)$ the three pentagon
arguments are
\[
xy=q^3,\qquad
\frac{x(1-y)}{1-xy}=q^2,\qquad
\frac{y(1-x)}{1-xy}=q^5.
\]
Hence $R_1+R_2=R_3+R_2+R_5$, or $R_1=R_3+R_5$.
Reflection gives $R_1+R_5=R_2+R_3=\zeta(2)$, so
$R_2=2R_5$.  Expanding the logarithmic corrections by
\eqref{eq:plastic-complements} gives
\[
\operatorname{Li}_2(q^2)+3\lambda^2
=2\operatorname{Li}_2(q^5)+5\lambda^2,
\]
as required.
\end{proof}

\subsection{The trilogarithm formula: one Kummer specialization}

For real nonzero $z\ne1$, define the single-valued trilogarithm
\[
P(z)=\operatorname{Re}\operatorname{Li}_3(z)
 -\log|z|\operatorname{Re}\operatorname{Li}_2(z)
 -\frac13\log^2|z|\log|1-z|,
\]
with the continuous value $P(1)=\zeta(3)$.
The inversion and duplication equations are
\[
P(z^{-1})=P(z),\qquad P(z)+P(-z)=\frac14P(z^2).
\]
Kummer's nine-term equation, in this single-valued real form, is
\begin{align}
&2P(x)+2P(y)
 +2P\!\left(\frac{x(1-y)}{x-1}\right)
 +2P\!\left(\frac{y(1-x)}{y-1}\right)\notag\\
&\quad+2P\!\left(\frac{1-x}{1-y}\right)
 +2P\!\left(\frac{x(1-y)}{y(1-x)}\right)
 -P(xy)-P(x/y)\notag\\
&\quad-P\!\left(\frac{x(1-y)^2}{y(1-x)^2}\right)
 =2\zeta(3).
\label{eq:plastic-Kummer}
\end{align}
We use the form recorded by Zagier, Section~7, Example~3,
equation~(7) \cite{ZagierPolylogs1991}.

The single-valued Landen equation is
\[
P(x)+P(1-x)+P(1-x^{-1})=\zeta(3),\qquad0<x<1.
\]
Its ordinary real counterpart is
\begin{align}
\operatorname{Li}_3(x)+\operatorname{Li}_3(1-x)
 +\operatorname{Li}_3(1-x^{-1})
 &=\zeta(3)+\zeta(2)\log x\notag\\
 &\quad-\frac12\log^2x\log(1-x)+\frac16\log^3x.
\label{eq:plastic-ordinary-Landen}
\end{align}
One obtains it by differentiation using dilogarithm reflection and
Landen, with the integration constant determined at $x=1$.
See also Nakamura--Shiraishi \cite{NakamuraShiraishi2025}.

\begin{theorem}[The manuscript's plastic trilogarithm formula]
\label{thm:plastic-trilog}
\begin{equation}
-\operatorname{Li}_3(q)+\frac32\operatorname{Li}_3(q^2)
 +\frac13\operatorname{Li}_3(q^3)-\operatorname{Li}_3(q^5)
=\lambda^3.
\label{eq:plastic-trilog-proved}
\end{equation}
\end{theorem}

\begin{proof}
Put $P_k=P(q^k)$.  At $(x,y)=(q,q^2)$, the nine arguments in
\eqref{eq:plastic-Kummer}, in displayed order, are
\[
q,\quad q^2,\quad -q^{-1},\quad -q^4,\quad q^2,
\quad q^{-3},\quad q^3,\quad q^{-1},\quad q^{-5}.
\]
Using inversion and duplication, the equation becomes
\[
\mathsf K=-P_1+\frac92P_2+P_3-2P_4-P_5+\frac12P_8
=2\zeta(3).
\]
Landen at $x=q$, with $1-q=q^5$, similarly gives
\[
\mathsf A=P_1+P_5-P_4+\frac14P_8=\zeta(3).
\]
Consequently the exact row operation $(\mathsf K-2\mathsf A)/3$
proves
\begin{equation}
-P_1+\frac32P_2+\frac13P_3-P_5=0.
\label{eq:plastic-P3-proved}
\end{equation}

It remains to convert this to ordinary trilogarithms.  Denote the
left side of \eqref{eq:plastic-trilog-proved} by $D$, and put
\[
S=-\operatorname{Li}_2(q)+3\operatorname{Li}_2(q^2)
 +\operatorname{Li}_2(q^3)-5\operatorname{Li}_2(q^5).
\]
The two ordinary reflection equations and
Theorem~\ref{thm:plastic-dilog} give
\begin{align*}
\operatorname{Li}_2(q)+\operatorname{Li}_2(q^5)
 &=\zeta(2)-5\lambda^2,\\
\operatorname{Li}_2(q^2)+\operatorname{Li}_2(q^3)
 &=\zeta(2)-6\lambda^2,\\
\tfrac12\operatorname{Li}_2(q^2)-\operatorname{Li}_2(q^5)
 &=\lambda^2.
\end{align*}
Taking minus the first row, plus the second, plus four times the
third yields $S=3\lambda^2$.
The remaining logarithmic coefficient in the definition of $P$ is
\[
T=\{-1\cdot1^2\cdot5
 +(3/2)\cdot2^2\cdot3
 +(1/3)\cdot3^2\cdot2
 -1\cdot5^2\cdot1\}\lambda=-6\lambda.
\]
Thus the left side of \eqref{eq:plastic-P3-proved} is
\[
D-\lambda S-\frac{\lambda^2}{3}T
=D-3\lambda^3+2\lambda^3=D-\lambda^3.
\]
This proves \eqref{eq:plastic-trilog-proved}.
\end{proof}

The proof promotes the two displayed manuscript formulas from
numerical support to exact consequences of classical functional
equations.  It does not rely on pseudointegration at an isolated
algebraic base.

\subsection{Two log-free identities}

\begin{theorem}[Log-free plastic ladders]
\label{thm:plastic-logfree}
At the same base $q$, the following identities hold:
\begin{align}
&6\operatorname{Li}_2(q)-5\operatorname{Li}_2(q^2)
 -5\operatorname{Li}_2(q^3)+6\operatorname{Li}_2(q^5)
 =\zeta(2),\label{eq:plastic-logfree2}\\
&12\operatorname{Li}_3(q)-5\operatorname{Li}_3(q^2)
 -4\operatorname{Li}_3(q^3)-8\operatorname{Li}_3(q^4)
 +8\operatorname{Li}_3(q^5)+2\operatorname{Li}_3(q^8)
 =4\zeta(3).
\label{eq:plastic-logfree3}
\end{align}
\end{theorem}

\begin{proof}
Six times the first ordinary reflection equation minus five times
the second proves \eqref{eq:plastic-logfree2}.
For the trilogarithm, write $L_k=\operatorname{Li}_3(q^k)$.
Equation~\eqref{eq:plastic-ordinary-Landen} at $x=q$ and $x=q^2$,
with duplication to eliminate negative arguments, gives
\begin{align*}
A&=L_1+L_5-L_4+\frac14L_8
  =\zeta(3)+\zeta(2)\lambda-\frac73\lambda^3,\\
B&=\frac54L_2+L_3-L_1
  =\zeta(3)+2\zeta(2)\lambda-\frac{14}3\lambda^3.
\end{align*}
The logarithmic tails are proportional.  Therefore
$4(2A-B)=4\zeta(3)$, exactly the claimed six-argument identity.
\end{proof}

These formulas are explicit additional identities for the package.
Their derivations and the classical origins of the functional
equations are given in full; no global priority claim is made.
They also give exponentially convergent algebraic series, for example
\begin{equation}
\zeta(3)=\sum_{n=1}^{\infty}
\frac{3q^n-\tfrac54q^{2n}-q^{3n}-2q^{4n}
                 +2q^{5n}+\tfrac12q^{8n}}{n^3}.
\label{eq:plastic-zeta3-series}
\end{equation}
Absolute convergence justifies expansion of the polylogarithms.
If this series is truncated after $N$ terms, its absolute tail is
at most
\[
\frac{39}{4}\,
\frac{q^{N+1}}{(N+1)^3(1-q)},
\]
obtained by summing the absolute coefficients and using
$q^{kn}\le q^n$ for $k\ge1$.  The constant $39/4$ is deliberately
simple; the sharper sum of the six individual geometric tails can
also be used.

\subsection{Exact certificates and their scope}

The accompanying script checks the polynomial identities and all
three Rogers and nine Kummer argument reductions in
$\mathbb Q[X]/(X^3+X^2-1)$.  It constructs the Kummer row from the
signed exponent list and performs the displayed rational row
operations.  Separately, it enumerates a bounded collection of
trinomials, checks the predicted cyclotomic factors and irreducible
quotients with exact polynomial arithmetic, and groups exact
minimal polynomials to verify the expected collision table.
The finite enumeration is a reproducible consistency check; the
unbounded completeness result is Theorem~\ref{thm:complement-collisions}.
